% =====================================================================
% sections/literature.tex — long version only (~3-4 pages)
% =====================================================================

\section{Related literature}
\label{sec:literature}

\subsection{Early-life shocks and adult human capital}

The conceptual basis for studying the long-run effects of early-life
exposure to a public-health intervention is the fetal-origins
literature initiated by \textcite{barker1995}. Barker's evidence from
British cohort data showed that nutritional deprivation in utero was
associated with elevated risks of coronary heart disease and
hypertension decades later. The economics literature has extended this
``fetal origins'' framing to a range of early-life environmental
shocks. \textcite{maccini2009} use Indonesian rainfall variation in
the year of birth as a quasi-experimental shock to early-life
nutrition; they find that women born in years of higher rainfall---and
hence better childhood nutrition---attain more years of schooling and
better socioeconomic status as adults. \textcite{majid2015} uses
Ramadan timing as a source of in-utero exposure to maternal fasting
and finds that prenatal exposure reduces adult labor supply and
shifts the composition of work. \textcite{duflo2001}'s study of
Indonesia's INPRES primary-school construction program is a closely
related design that uses geographic and cohort variation to identify
the effect of a public-good expansion on adult outcomes; she finds
substantial returns in years of schooling and wages. The Village
Midwife Program is a positive shock to the early-life environment of
the same kind that this literature has studied, with the difference
that the policy variation is plausibly larger and more sustained
than rainfall or Ramadan timing.

\subsection{The Village Midwife Program in particular}

A short empirical literature has examined the Village Midwife Program
directly. \textcite{frankenberg2005}, using IFLS waves 1 and 2 and
focusing on the first ten years of the program, document a positive
effect of midwife arrival on children's height-for-age, with the
effect concentrated among children whose mothers had access to the
midwife during pregnancy. \textcite{thomas2001} reports that mothers'
body mass index rose in communities that received midwives, suggesting
the nutrition-counseling component had measurable effects on adult
women as well as children. \textcite{triyana2016} returns to the same
question ten years later, exploiting the additional variation
introduced by the Asian Financial Crisis-induced contraction in
midwife stocks, and finds that the program's effect on antenatal-care
use persists at the ten-year mark but most adult-health effects do
not---a finding that is partly attributable to the contraction itself.

\textcite{ahsan2021} is the most recent paper in this lineage and the
direct predecessor to the present analysis. They estimate a single
specification (\texttt{areg} with community absorbed) on the cohorts
born 1989--1996 and report that in-utero exposure to the village
midwife produced long-run improvements in adult health, cognition,
and economic outcomes among men. They do not report sibling fixed
effects, do not implement a heterogeneity-robust staggered-DiD
estimator, and do not study mental-health or personality outcomes. The
present analysis revisits their findings with a layered identification
stack, an expanded outcome set, and a transparent treatment of
identification weaknesses.

\subsection{Mental-health and socioemotional development}

A more recent literature has begun to document the long-run
consequences of early-life environments for mental-health and
socioemotional outcomes, complementing the older focus on physical
health and cognition. \textcite{casepaxson2008} document that
adult height---a marker of early-life nutritional environment---is
correlated with adult labor-market outcomes through cognitive channels,
suggesting that early-life shocks operate through pathways beyond
those captured by adult height itself. \textcite{premand2022} report
experimental evidence from a low-income setting that combining
behavioral-change promotion with cash transfers produces measurable
gains in early-childhood development that include socioemotional
outcomes. \textcite{sorrenti2025} provide quasi-experimental evidence
that explicit socioemotional-skills training in early adolescence
generates lasting effects on educational success.
\textcite{webb2024} uses weather shocks during early life in
Indonesia---an identification strategy similar to
\textcite{maccini2009}---to document critical periods in cognitive and
socioemotional development; the present paper's choice to standardize
the depression and Big-Five outcomes in the IFLS adult module follows
his measurement approach. The HeadStart program documentation
\autocite{headstart} provides a synthesis of the developmental case
for treating socioemotional outcomes as a first-order metric of
early-life intervention success.

The case for adding mental-health and personality outcomes to the
study of the Village Midwife Program rests on three observations.
First, the program's service bundle---nutrition counseling,
infection control, prenatal monitoring---includes the same factors
that the developmental-medicine literature has linked to long-run
mental-health outcomes. Second, the IFLS adult module from wave~5
contains validated psychometric instruments (CESD-10 for depression,
Big-Five Inventory for personality) that can be standardized and
compared across cohorts. Third, the conventional outcomes used in
\textcite{ahsan2021}---adult height, cognition, labor-market
status---may understate the program's true human-capital effects if
the bundle's largest impact is on socioemotional rather than physical
or cognitive margins.

\subsection{Staggered difference-in-differences methodology}

The recent econometrics literature on staggered-treatment
difference-in-differences forms the methodological backbone of the
present paper. \textcite{goodmanbacon2021} demonstrates that the
canonical TWFE estimator with staggered treatment timing is biased in
the presence of heterogeneous treatment effects, because the
estimator implicitly uses already-treated units as controls for
later-treated ones. The bias can be severe enough to flip the sign of
the estimated effect.

A wave of alternative estimators has followed. The
\textcite{callawaysantanna2021} approach computes group-time average
treatment effects $ATT(g,t)$ using a clean comparison group (either
never-treated or not-yet-treated) and aggregates them with explicit
weights. The \textcite{dechaisemartin2021} estimator constructs
analogous unbiased aggregations under different assumptions about
treatment-effect dynamics. The Sun--Abraham aggregation provides a
cohort-aggregated event-study version of the same idea. All three
appear in this paper as robustness checks against the headline TWFE
estimate.

For the multiple-outcomes problem within a domain---several physical
health measures collectively comprising ``health,'' several depression
items collectively comprising ``mental health''---I follow
\textcite{anderson2008} in using inverse-covariance-weighted
composite indices. This approach is now standard in early-life
intervention studies and has the dual advantage of reducing
multiple-testing inflation and producing scale-free composite outcomes
that are interpretable across studies.
